\documentclass[12pt,a4paper,oneside]{article}
\usepackage[brazil]{babel}
\usepackage[utf8]{inputenc}
\usepackage{olymp}

\setcounter{problem}{2}

\begin{document}

\begin{problem}{Futebol olímpico}{futebol}{2}
 
O Brasil é uma potência mundial no futebol: o único a participar de todas as Copas do Mundo e o único pentacampeão. Em Olimpíadas, a história era bem diferente... Até que, finalmente, o Brasil ganhou sua primeira medalha de ouro nesta modalidade. E foi em casa, na Rio 2016! E não parou aí, ganhou outra logo em seguida, na Tokyo 2020, disputada em 2021.

O futebol é uma modalidade olímpica com limitação de idade. Os times devem ser formados por jogadores com menos de 23 anos. A partir de Atlanta 96, permite-se que cada time tenha até 3 atletas com 23 anos ou mais. Em outras palavras, um time não é aceito se o número de jogadores do time que já completaram 23 anos for maior que 3.

Faça um programa para testar esta regra com diferentes limites de diferentes idades.
%\Notes
%
%Observações, se tiver.

\InputFile

A entrada contém duas linhas de números inteiros informando os dados de um time. A primeira linha contém três números, $N$, $I$ e $L$, indicando respectivamente o número de jogadores do time, a idade usada na regra e o limite de jogadores que podem ter completado esta idade. A segunda linha contém uma lista de $N$ números inteiros, que são as idades dos jogadores do time. Restrições: $1 \leq N \leq 30, 15 \leq I \leq 30, 0 \leq L \leq N$.


\OutputFile

Seu programa deve verificar se entre os $N$ jogadores existem mais que $L$ %jogadores
que já completaram $I$ anos %. Seu programa deve 
e
escrever ``\texttt{Aceito.}'' ou ``\texttt{Nao aceito.}'', indicando se o time atende ou não a regra.

\Notes
Note que a mensagem escrita não tem acento e tem um ponto final.

\Example

\begin{example}
\exmp{
6 21 2
15 23 22 18 19 21
}{
Nao aceito.
}%
\end{example}

\begin{example}
\exmp{
6 22 2
15 23 22 18 19 21
}{
Aceito.
}%
\end{example}

\begin{example}
\exmp{
6 20 3
15 23 22 18 19 21
}{
Aceito.
}%
\end{example}

\begin{example}
\exmp{
8 15 1
14 14 14 14 14 14 14 15
}{
Aceito.
}%
\end{example}

\begin{example}
\exmp{
5 15 1
8 8 15 20 8
}{
Nao aceito.
}%
\end{example}



%Comentários sobre o exemplo, se necessário.

\small
No primeiro examplo, o time não é aceito porque há mais de 2 jogadores que completaram 21 anos.

\end{problem}

\end{document}
